\documentclass[aspectratio=169]{beamer}
\usepackage{animate}

% Show notes
\usepackage{pgfpages}
%\setbeameroption{show notes on second screen}

\usepackage[utf8]{inputenc}

\setbeamertemplate{note page}
{%
  % Header box
  \insertvrule{.27\paperheight}{white!90!black}
  \vskip-0.286\paperheight
  \begin{minipage}{0.4\textwidth}
    \insertslideintonotes{0.25}%
  \end{minipage}
  \begin{minipage}{0.55\textwidth}
    {\textbf \insertshortframetitle}
  \end{minipage}

  % Title
  \begin{minipage}{0.88\paperwidth}
  \insertnote
  \end{minipage}
}

\newcommand{\timenote}[3]{
    \note{
        \vskip-0.1\paperheight
        \hspace{4.44cm}Time: #1 -- #2 min
        \vskip0.086\paperheight
        {\scriptsize #3}}
}

\usepackage{mathtools}


\usepackage[]{algorithm2e}
\usepackage{bm}
\renewcommand{\vec}[1]{%
  \ifcat\noexpand#1\relax % check if the argument is a control sequence
    \bm{#1} % probably Greek
  \else
    \mathbf{#1} % single character
  \fi
}
\usepackage{amssymb}
\usepackage{amstext}
\usepackage{amsmath}
\usepackage{amsthm}
\usepackage{changepage}
\usepackage{mathrsfs}  % Used for mathscr 'L'
\usepackage{gensymb}  % degree sign
\usepackage{bm}
\usepackage{bbm}
\renewcommand{\epsilon}{\varepsilon}
\usepackage{cancel}
%\usepackage{breqn}
%\usepackage{enumitem}
\usepackage{booktabs}
%\usepackage[table]{xcolor}
\usepackage{multirow}

\usepackage{empheq}
%\usepackage[most]{tcolorbox}
%\newtcbox{\mymath}[1][]{%
%    nobeforeafter, math upper, tcbox raise base,
%    enhanced, colframe=blue!30!black,
%    colback=blue!30, boxrule=1pt,
%    #1}
%\tcbset{
%    enhanced,
%    breakable,
%    attach boxed title to top center,
%    top=1mm,
%    coltitle=black,
%    beforeafter skip=\baselineskip,
%}

\newcommand{\etal}{\emph{et~al.}}

% Graphics
\graphicspath{{../images/}{../matlab/}{../../camera_ready/images/}}
\DeclareGraphicsExtensions{.pdf,.jpeg,.png,.eps}

\newcommand{\FULLSCREEN}{0.65\textwidth}

\makeatletter
\newcommand{\vast}{\bBigg@{4}}
\newcommand{\Vast}{\bBigg@{5}}
\newcommand{\VAST}{\bBigg@{16}}
\newcommand{\VAst}{\bBigg@{14}}
\makeatother

\usepackage[export]{adjustbox}[2011/08/13]

\newcommand{\cammat}[2]{\left[#1\;|\;#2\right]}
\newcommand{\R}{\ensuremath{\mathbbm{R}}}
\newcommand{\K}{\ensuremath{\mathbbm{K}}}
\newcommand{\C}{\ensuremath{\mathbbm{C}}}
\renewcommand{\phi}{\varphi}
\DeclareMathOperator*{\argmax}{arg\,max}
\DeclareMathOperator*{\argmin}{arg\,min}
\DeclareMathOperator{\nulldim}{nulldim}
\DeclareMathOperator{\adj}{adj}
\newcommand{\mat}[1]{\bm{#1}}
\newcommand{\nullspace}[1]{\mathcal{N}{#1}}
\newcommand{\lr}{\ensuremath{\mathrm{\quad\Leftrightarrow\quad}}}
\newcommand{\A}{\ensuremath{\mathcal{A}}}
\newcommand{\N}{\ensuremath{\mathcal{N}}}
\newcommand{\reg}{\ensuremath{\mathcal{R}}}
\newcommand{\Real}{\ensuremath{\mathbbm{R}}}
\newcommand{\Z}{\ensuremath{\mathbbm{Z}}}
\newcommand{\Tp}{T}
\newcommand{\norm}[1]{\|#1\|}
\newcommand{\frobnorm}[1]{\norm{#1}_F}
\newcommand{\fr}[2]{\frac{#1}{#2}}
\newcommand{\pf}[2]{\frac{\partial #1}{\partial #2}}
\DeclareMathOperator{\tr}{tr}
\DeclareMathOperator{\diag}{diag}
\DeclareMathOperator{\atan}{atan2}
\DeclareMathOperator{\rank}{rank}

\newcommand{\T}{T}
\usepackage{setspace}
%\usepackage{xcolor}
\usepackage{tikz}
\usepackage{pgfplots}

%\usepackage{caption}
\usepackage{subcaption}

\def\thicknesslineplot{thick}
\def\thicknesshack{0.085cm}
\def\mythickness{ultra thick}
\def\thicknesslineplotTwo{thick}
\def\mythicknesssmall{thick}
\def\thicknesslineplotThree{very thick}
\def\boxplotheight{8.5cm}
\def\boxplotwidth{7.5cm}
\def\boxplotheightreal{3.5cm}

\definecolor{gtcolor}{rgb}{0,0,0}
\definecolor{ourcolor}{rgb}{0,0.4,0.6}
\definecolor{fitzgibboncolor}{rgb}{0,0.7,0.7}
\definecolor{dingcolor}{rgb}{1,0.2,0.3}
\definecolor{kukelovacolor}{rgb}{0,0.6,0}
\definecolor{bougnouxcolor}{rgb}{0,0.7,0.7}
\definecolor{valtonenoernhagcolor}{rgb}{0.7,0,0.7}
\definecolor{inliercolor}{rgb}{0.2,1,0.1}
\definecolor{outliercolor}{rgb}{1,0,0}
\definecolor{smallcolor1}{rgb}{0.,0,0}
\definecolor{smallcolor2}{rgb}{0,0,0}
\definecolor{hiddencolor}{rgb}{0.1,0.9,0.1}
\definecolor{elimcolor}{rgb}{0.6,0.1,0.6}

\makeatletter
\DeclareRobustCommand\eg{\emph{e.g}\@ifnextchar.{}{.\@}}
\DeclareRobustCommand\etal{\emph{et~al}\@ifnextchar.{}{.\@}}
\DeclareRobustCommand\ie{\emph{i.e}\@ifnextchar.{}{.\@}}
\DeclareRobustCommand\cf{\emph{cf}\@ifnextchar.{}{.\@}}
\DeclareRobustCommand\NB{\emph{N.B}\@ifnextchar.{}{.\@}}
\DeclareRobustCommand\wrt{w.r.t\@ifnextchar.{}{.\@}}
\makeatother

\newcommand{\CE}{\ensuremath{\mat{C}_\mathcal{E}}}
\newcommand{\CR}{\ensuremath{\mat{C}_\mathcal{R}}}
\newcommand{\CB}{\ensuremath{\mat{C}_\mathcal{B}}}
\newcommand{\CEk}[1]{\ensuremath{\mat{C}_{\mathcal{E}_{#1}}}}
\newcommand{\CRk}[1]{\ensuremath{\mat{C}_{\mathcal{R}_{#1}}}}
\newcommand{\CBk}[1]{\ensuremath{\mat{C}_{\mathcal{B}_{#1}}}}
\newcommand{\XE}{\ensuremath{\mat{X}_\mathcal{E}}}
\newcommand{\XR}{\ensuremath{\mat{X}_\mathcal{R}}}
\newcommand{\XB}{\ensuremath{\mat{X}_\mathcal{B}}}

\newcommand{\sing}[1]{\vec{\sigma}{#1}}
\usepackage{standalone}
\pgfplotsset{
    table/search path={../../graphs},
}

\makeatletter
\@namedef{ver@everyshi.sty}{} %Prevents some error caused by tikz, has to be agter xcolor
\makeatother
\usepackage{tikz}  % Used for problem geometry picture
\usepackage{pgfplots}
\usepackage{tikzpagenodes}
\usetikzlibrary{positioning}
\usetikzlibrary{arrows,shapes}
\usetikzlibrary{arrows.meta}
\tikzset{%
  thick arrow/.style={
     -{Triangle[angle=90:2pt 1.5]},
     line width=4mm,
     draw=gray,
     draw opacity=0.7,
  },
}
\pgfplotsset{compat=1.15}
\usetikzlibrary{shapes.misc, positioning}
\usetikzlibrary{arrows}
\usetikzlibrary{calc}

\pgfdeclarelayer{foreground}  % Used for correct opacity
\pgfdeclarelayer{background}  % Used for correct opacity
\pgfsetlayers{background,main,foreground}

\DeclareGraphicsRule{.tif}{png}{.png}{`convert #1 `dirname #1`/`basename #1 .tif`.png}

\usepackage{environ}  % Used for scaletikzpicturewidth
\makeatletter
\newsavebox{\measure@tikzpicture}
\NewEnviron{scaletikzpicturetowidth}[1]{%
  \def\tikz@width{#1}%
  \def\tikzscale{1}\begin{lrbox}{\measure@tikzpicture}%
  \BODY
  \end{lrbox}%
  \pgfmathparse{#1/\wd\measure@tikzpicture}%
  \edef\tikzscale{\pgfmathresult}%
  \BODY
}
\makeatother

\tikzset{add reference/.style={insert path={%
coordinate [pos=0,xshift=-0.5\pgflinewidth,yshift=-0.5\pgflinewidth] (#1 south west)
coordinate [pos=1,xshift=0.5\pgflinewidth,yshift=0.5\pgflinewidth]   (#1 north east)
coordinate [pos=.5] (#1 center)
(#1 south west |- #1 north east)     coordinate (#1 north west)
(#1 center     |- #1 north east)     coordinate (#1 north)
(#1 center     |- #1 south west)     coordinate (#1 south)
(#1 south west -| #1 north east)     coordinate (#1 south east)
(#1 center     -| #1 south west)     coordinate (#1 west)
(#1 center     -| #1 north east)     coordinate (#1 east)
}}}

\usepackage{listings}
\usepackage{multimedia}
\usepackage{graphicx}

\usepackage{colortbl}
\definecolor{cA}{RGB}{212, 237, 244}%
\definecolor{cAa}{RGB}{192, 217, 224}%
\definecolor{cB}{RGB}{237, 212, 244}%
\definecolor{cC}{RGB}{244, 237, 212}%
\definecolor{darkgreen}{RGB}{0, 170, 0}
\newcommand\colorA{\cellcolor{cA}}
\newcommand\colorAa{\cellcolor{cAa}}
\newcommand\colorB{\cellcolor{cB}}
\newcommand\colorC{\cellcolor{cC}}
\newcommand{\ai}{{\color{darkgreen}a_i}}
\newcommand{\Twoai}{{\color{darkgreen}2a_i}}
\newcommand{\bi}{{\color{orange}b_i}}

\DeclareMathOperator{\syz}{Syz}
\DeclareMathOperator{\lt}{LT}
\DeclareMathOperator{\sat}{Sat}

% Matlab code
\usepackage[numbered,framed]{matlab-prettifier}
\lstset{
style              = Matlab-editor,
%  basicstyle         = \mlttfamily,
%  escapechar         = ",
%  mlshowsectionrules = true,
}

\usepackage{tikz-3dplot}
\usetikzlibrary{calc,arrows.meta,positioning,backgrounds}
\pgfdeclarelayer{foreground}  % Used for correct opacity
\pgfdeclarelayer{background}  % Used for correct opacity
\pgfsetlayers{background,main,foreground}

\usetheme{metropolis}           % Use metropolis theme
\title{\vspace{0.3cm}{\small Practical aspects of solving a system of polynomial equations}\vspace{-0.5cm}}
\author{\vspace{-0.3cm}Marcus Valtonen Örnhag}
\institute{%
    Centre for Mathematical Sciences, Lund University}
\date{\vspace{0.2cm}{\footnotesize\emph{Workshop in Algebraic Geometry, May 2021}}}
\titlegraphic{%
    \includegraphics[height=1.7cm]{LundUniversity_C2line_CMYK}
}

\begin{document}
\maketitle

\section{Introduction}

\begin{frame}{Introduction}
Today's tutorial/workshop will contain:
\begin{itemize}
\item A brief theoretical walkthrough,
\item Tips and tricks on how to practically speed-up your solver,
\item Hands-on experience with the automatic Gröbner basis generator,
\item A group exercise followed by a discussion session.
\end{itemize}
\end{frame}

\section{Theory}
\begin{frame}{Recap: Polynomial systems of equations}
Using multi-index notation,
let~$\mat{x^{\alpha}}$ denote a \emph{monomial} of degree $|\mat{\alpha}|$.
Our main goal is to solve a polynomial system of equations
\begin{equation}\label{kappa:eq:polsys}
    \begin{aligned}
    f_1(\mat{x}) & = 0, \\
                 & \hspace{0.575em} \vdots \\
    f_s(\mat{x}) & = 0,
    \end{aligned}
\end{equation}
where each polynomial equation can be expressed
as~$f=\sum_{\mat{\alpha}}c_{\mat{\alpha}}\mat{x^\alpha}$.

\end{frame}

\begin{frame}{Recap: The toolbox}
\begin{itemize}
\item  The set of all solutions to~\eqref{kappa:eq:polsys} is called an \emph{affine variety}, and is denoted
$\mat{V}(f_1,\ldots,f_s)\subset\C$.
\item An \emph{ideal} $I\subset\C[\mat{x}]$ is an additive group satisfying the \emph{absorption property},~\ie{}~if~$f\in I$ and $h\in\C[\mat{x}]$ then $hf\in I$. Moreover, given a set of polynomials $f_1,\ldots,f_s\in\C[\mat{x}]$
we consider the ideal
\begin{equation*}
\langle f_1,\ldots,f_s\rangle = \left\{\sum_{i=1}^sh_if_i\;:\;h_1,\ldots,h_s\in\C[\mat{x}]\right\},
\end{equation*}
which we will refer to as the \emph{ideal generated by} $f_1,\ldots,f_s$.

\end{itemize}
Here~$\C[\mat{x}]$ denotes the set
of polynomials in~$\mat{x}$ with coefficients in $\C$.
\end{frame}


\begin{frame}{Recap: The toolbox (cont.)}

\begin{itemize}
\item The \emph{Hilbert Basis Theorem} states that every ideal of $\C[\mat{x}]$ is
finitely generated, \ie{} given an ideal~$I\subset\C[\mat{x}]$ there exist
$f_1,\ldots,f_s\in\C[\mat{x}]$ such that $I=\langle f_1,\ldots,f_s\rangle$.
\item The polynomial system of equations~\eqref{kappa:eq:polsys} is defined by the
generated ideal; however, the basis is in general not unique.
\end{itemize}

\alert{Solution}: Impose a \emph{monomial ordering} $\longrightarrow$ unique representation.
\end{frame}

\begin{frame}{Recap: The toolbox (cont.)}

\begin{itemize}
\item The \emph{leading term} (\wrt{} the monomial ordering) is well-defined. We denote it~$\lt(\cdot)$.
\item For an ideal $I$ define $\lt(I)\coloneqq\{\lt(f)\;:\;f\in I\}$.
\item A \emph{Gröbner basis} is a finite subset \mbox{$\mat{G}=\{g_1,\ldots,g_t\}\subset I$}
satisfying
\begin{equation*}
    \langle\lt(g_1),\ldots,\lt(g_t)\rangle=\lt(I)\;.
\end{equation*}
\end{itemize}
\end{frame}

\begin{frame}{Recap: Why is all this necessary?}
How does this help us solve the original system?

$\quad\longrightarrow\quad$ We can compute a Gröbner basis efficiently using~\emph{Buchberger's algorithm}.
\end{frame}

\begin{frame}{Recap: The action matric method}
Assume zero-dimensional ideals, \ie{} finitely many solutions to the original system.
\begin{itemize}
\item Define the \emph{coset} $[f] \coloneqq \{f+h\;|\;h\in I\}$.
\item The linear operator $T_f\::\:\C[\mat{x}]/I\longrightarrow\C[\mat{x}]/I$, defined as $T_g([g])=[fg]$,
satisfies $T_f=T_g$ iff $f-g\in I$.
% \item The quotient space $\C[\mat{x}]/I$ is finite dimensional if the original system of
% equations have finitely many solutions, hence
\item The operation $T_\alpha(\cdot)$ can be
represented using a matrix~$\mat{M}_\alpha$. This is the \emph{action matrix}.
\end{itemize}
\end{frame}

\begin{frame}{Recap: The action matric method (cont.)}
Select a monomial basis $\mathcal{B}=\{[\mat{x}^{\beta_i}]\}_{i=1}^K$ for~$\C[\mat{x}] / I$.
Then
\begin{equation*}
    T_\alpha([\mat{x}^{\beta_i}]) = [\alpha(\mat{x})\mat{x}^{\beta_i}] = \sum_{j=1}^Km_{ij}[\mat{x}^{\beta_j}]\;.
\end{equation*}
This implies that, for $\mat{x}\in V(I)$,
\begin{equation*}\label{kappa:eq:action}
    \alpha(\mat{x})\mat{x}^{\beta_i}=\sum_{j=1}^Km_{ij}\mat{x}^{\beta_j}\;.
\end{equation*}
Let $\mat{b}$ represent the basis~$\mathcal{B}$ (as a vector), then we have an \alert{eigenvalue problem}
\begin{equation*}\label{eq:actionmat}
    \alpha(\mat{x})\mat{b}(\mat{x}) = \mat{M}_\alpha\mat{b}(\mat{x})\;.
\end{equation*}
\end{frame}

\begin{frame}{The reality}
For any practical applications in computer vision, robotics and neighboring fields:
\begin{itemize}
\item We need fast or even real-time performance,
\item The input is contaminated with noise. $\longrightarrow$ Several different inputs are
required for robustness.
\end{itemize}

Unfortunately, \emph{Buchberger's algorithm} is NP-complete and often quite heavy. More
efficient algorithms exist, such as Faugère's F4 and F5 algorithms, but they are
still not feasible.
\end{frame}


\begin{frame}{The reality (cont.)}
But there are good news!
We are often interested in solving generic problem instances, \eg{} the relative pose
problem
\begin{equation*}
\begin{aligned}
    \mat{x}'_i\mat{E}\mat{x}_i &= 0, &i=1,\ldots,5\\
    \det(\mat{E}) &= 0, &\\
    2\mat{EE}^\T\mat{E}- \tr(\mat{EE}^\T)\mat{E} &= 0, &
\end{aligned}
\end{equation*}
Only the input data $\mat{x}_i$ and $\mat{x}_i'$ change!
\end{frame}

\begin{frame}{The reality (cont.)}
All non-degenearte instances of the same problem share the same characteristics:
\begin{itemize}
\item The same number of solutions,
\item Gröbner basis the same (\wrt{} a monomial ordering, often GREVLEX)${}^*$,
\item Creating the action matrix involves the same kind of computations${}^*$.
\end{itemize}
${}^*$Only differs for each instance in terms of the coefficients from the input data.
\end{frame}

\begin{frame}{Elimination templates}
\alert{Idea:} We try to find the action matrix using \emph{elimination templates},
which encode the relevant eliminations directly.
\end{frame}

\begin{frame}{Example}
\end{frame}

\begin{frame}{Elimination templates (cont.)}
Consider the $i$:th row in $\alpha\mat{b} = \mat{M}_\alpha\mat{b}$.
Any way of expressing $[\alpha b_i]$ reduces to find a polynomial
\begin{equation}\label{eq:polys}
    p_i = \alpha b_i -\sum_{j=1}^Km_{ij}b_{j}\in I\;.
\end{equation}
\alert{Goal:} Find a \emph{single} polynomial like this expressed in elements from
$\mat{b}$ of $\alpha\mat{b}$.
\end{frame}

\begin{frame}{Elimination templates (cont.)}
\begin{columns}
    \begin{column}{0.5\textwidth}
\alert{Solution:}
Augment the original system
for some monomials~$m_1,\ldots, m_\ell$.
    \end{column}%
    \begin{column}{0.5\textwidth}
\footnotesize
\begin{equation*}\label{kappa:eq:polsys-aug}
    \begin{aligned}
    \left.
    \begin{aligned}
    f_1(\mat{x}) & = 0, \\
                 & \hspace{0.575em} \vdots \\
    f_s(\mat{x}) & = 0, \\
    \end{aligned}
    \right\}&\text{original}
    \\
    \left.
    \begin{aligned}
    m_1f_1(\mat{x}) & = 0, \\
                 & \hspace{0.575em} \vdots \\
    m_1f_s(\mat{x}) & = 0, \\
                 & \hspace{0.575em} \vdots \\
    m_\ell f_1(\mat{x}) & = 0, \\
                 & \hspace{0.575em} \vdots \\
    m_\ell f_s(\mat{x}) & = 0,
    \end{aligned}
    \right\}&\text{augmented}
    \end{aligned}
\end{equation*}
    \end{column}
\end{columns}
\end{frame}

\begin{frame}{Elimination templates (cont.)}
Let the augmented system be represented as~$\mat{CX}=0$ and partition it as follows
\begin{equation*}
    \begin{bmatrix}
        \CE &
        \CR &
        \CB
    \end{bmatrix}
    \begin{bmatrix}
        \XE \\
        \XR \\
        \XB
    \end{bmatrix}
    = 0,
\end{equation*}
where
\begin{itemize}
\item $\XB$ are the basis monomials,
\item $\XR$ are the \emph{reducible} monomials,
\item $\XE$~\emph{excessive} monomials.
\end{itemize}
\end{frame}

\begin{frame}{Elimination templates (cont.)}
We don't want the excessive monomials $\longrightarrow$ eliminate them
\begin{equation*}
    \begin{bmatrix}
        \mat{U}_\mathcal{E}' &
        \CRk{1}' &
        \CBk{1}' \\
        \mat{0} &
        \CRk{2}' &
        \CBk{2}'
    \end{bmatrix}
    \begin{bmatrix}
        \XE \\
        \XR \\
        \XB
    \end{bmatrix}
    = 0\;.
\end{equation*}
We may now discard the first row.
\end{frame}

\begin{frame}{Elimination templates (cont.)}
Assuming $\CRk{2}'$ is of full rank
\begin{equation*}
    \begin{bmatrix}
        \mat{I} &
        \CRk{2}'^{-1}\CBk{2}'
    \end{bmatrix}
    \begin{bmatrix}
        \XR \\
        \XB
    \end{bmatrix}
    = 0,
\end{equation*}
or
\begin{equation*}
    \XR = -\CRk{2}'^{-1}\CBk{2}'\XB\;.
\end{equation*}
Now, the reducible monomials can be expressed as a linear combination of the basis elements,
which is what we were looking for.
\end{frame}

\begin{frame}{Example}
\end{frame}

\begin{frame}{Elimination templates (cont.)}
\begin{itemize}
\item Which monomials should I use to expand the system?
    \begin{itemize}
    \item The fewer the better---fewer operations and better numerics.
    \end{itemize}
\item How to pick these monomials optimally?
    \begin{itemize}
    \item Computationally intensive.
    \end{itemize}
\end{itemize}
\end{frame}

\begin{frame}{Syzygy-based reduction}
\begin{definition}
Let~$(f_1,\ldots,f_t)$ be an ordered $t$-tuple of elements
$f_i\in\K[x]$, then the \emph{(first) syzygy module} is
\begin{equation*}
\syz(f_1,\ldots,f_t) \coloneqq \left\{
    (s_1,\ldots,s_t)\;|\; \sum_{i=1}^t s_if_i=0, \;\mathrm{ and }\; s_i\in\K[x]
\right\}\;.
\end{equation*}
\end{definition}
The syzygy module encodes the ambiguity of polynomials in an ideal.
Let $p\in \langle f_1,\ldots f_t\rangle$, then
\begin{equation*}
    p = \sum_{i=1}^t h_if_i = \sum_{i=1}^t (h_i+s_i)f_i, \qquad (s_1,\ldots,s_t)\in\syz(f_1,\ldots,f_t)\;.
\end{equation*}
\end{frame}

\begin{frame}{Example}
See Viktor Larsson's PhD thesis Example 1.1 (p. 46--47).
\end{frame}

\begin{frame}{Saturation of an ideal}
\begin{definition}
The \emph{saturation} of an ideal~$I\subset\K[x]$ \wrt{} the polynomial $f_s\in\K[x]$ is
defined as
\begin{equation*}
    \sat(I,\, f_s) \coloneqq \{p\;|\;\exists N\geq 0,\, f_s^N p\in I  \}\;.
\end{equation*}
\end{definition}
\alert{Interpretation:} we can remove solutions for which~$f_s(x) =0$.
\end{frame}

\section{Applications}
\begin{frame}{The importance of parameterization}
\begin{columns}
    \begin{column}{0.5\textwidth}
        \textbf{Example:} Three views with known rotation axis.\\[8mm]
        \begin{equation*}
            \begin{aligned}
                \mat{P}_1 &= [\mat{I}\;|\;\mat{0}] \\
                \mat{P}_2 &= [\mat{R}(\theta_2,\,{\color{red}\mat{v}})\;|\;\mat{t}_2] \\
                \mat{P}_3 &= [\mat{R}(\theta_3,\,{\color{red}\mat{v}})\;|\;\mat{t}_3]
            \end{aligned}
        \end{equation*}\\[8mm]
        Minimal case: One point and two lines.
    \end{column}%
    \begin{column}{0.5\textwidth}
        \centering
        \includestandalone[width=\linewidth]{tikz/3view_same_rot}
    \end{column}
\end{columns}
\end{frame}

\begin{frame}{The importance of parameterization}
\begin{columns}
    \begin{column}{0.5\textwidth}
        Choose a coordinate system in which~$\mat{v}=(0,\,1,\,0)$.
    \end{column}%
    \begin{column}{0.5\textwidth}
        \centering
        \includestandalone[width=\linewidth]{tikz/3view_same_rot}
    \end{column}
\end{columns}

\end{frame}

\section{Group exercise}

\begin{frame}{Warmup}
\textbf{Relative pose with one calibrated camera and one with unknown focal length}\\[8mm]
The fundamental matrix is~$\mat{E=FK}$, where $\mat{K}=\diag(f,\,f,\,1)$, and
\begin{equation*}
\begin{aligned}
    \mat{x}'_i\mat{F}\mat{x}_i &= 0, &i=1,\ldots,6\\
    \det(\mat{E}) &= 0, &\\
    2\mat{EE}^\T\mat{E}- \tr(\mat{EE}^\T)\mat{E} &= 0, &
\end{aligned}
\end{equation*}
\end{frame}

\begin{frame}{Problem 1}
\begin{columns}
    \begin{column}{0.5\textwidth}
        \textbf{Three (calibrated) cameras on a line}\\[8mm]
        \begin{itemize}
            \item How to parameterize?
            \item Which are the minimal cases?
        \end{itemize}
    \vspace{8mm}
    \footnotesize{(You may choose how many degrees of freedom you have in the rotation)}
    \end{column}%
    \begin{column}{0.5\textwidth}
        \centering
        \includestandalone[width=\linewidth]{tikz/3view_line}
    \end{column}
\end{columns}
\end{frame}

\begin{frame}{Problem 2}
\textbf{Panoramic stitching with a smartphone}\\[2mm]
Assume only rotation $\longrightarrow$ a homography~$\mat{H}$. Smartphones have an
IMU, hence a common reference direction exists. Then
\begin{equation*}
    \mat{H} = \mat{K}_2\mat{R}_2^\T\mat{R}_y\mat{R}_1\mat{K}_1^{-1},
\end{equation*}
where $\mat{R}_1$ and $\mat{R}_2$ is obtained from the IMU data. $\mat{R}_y$ is
a rotaiton about the $y$-axis.
Decide which case you want to look at
\begin{itemize}
    \item Focal length is unknown but the same,
    \item Focal length is unknown and possibly varying,
    \item Focal length and radial distortion parameter are unknown but the same,
    \item Focal length and radial distortion parameter are unknown and possibly varying.
\end{itemize}
\end{frame}

\end{document}
